\chapter{Convolution Absorbing Boundaries for Explicit-Circuit Quantum Simulation of the Wave Equation}\label{chap:qawave}

\begin{center}
    A paper submitted to \emph{Quantum}.

    Hoang Anh Nguyen\footnote{Department of Geophysics, Colorado School of Mines, Golden, Colorado 80401, USA}$^,$\footnote{Primary researcher and author for correspondence}
\end{center}

\subsection{Abstract}

Explicit quantum circuits for the wave equation via Hamiltonian simulation
[Sato \emph{et al.}, Phys.\ Rev.\ Research \textbf{6}, 033246 (2024)] are
restricted to closed (Dirichlet/Neumann/periodic) domains, in which
outgoing waves reflect back into the computational region.  We add
absorbing boundaries of the convolution type---the
complex-frequency-shifted perfectly matched layer (CPML) implemented
through exponential-kernel memory variables---and make the resulting
non-unitary dynamics quantum-implementable by Schr\"odingerisation.  We
obtain three results.  (i)~Explicit gate-level circuits for the absorbing
evolution, obtained by extending the Bell-basis term-evolution circuit of
Sato \emph{et al.}\ to projector-valued operator strings and Trotterising
to second order, demonstrated end to end on 17--21 qubits with
$10^{-3}$--$10^{-2}$ accuracy against exact references.  (ii)~A structural
obstruction: the memory-form CPML generator has an irreducibly indefinite
Hermitian part, $\lamax(H_1)=\Theta(\sqrt{\sigma_{\max}c/h})$ under any
diagonal rescaling of the memory fields, which by the recovery threshold
of Jin, Liu, and Ma costs $e^{2\lamax T}$ in post-selection.  (iii)~A
remedy: a classically precomputed Lyapunov symmetrizer renders the
transformed generator dissipative, replacing the $T$-dependent penalty by
a $T$-independent conditioning factor, measured at
$\kappa_2(S)=4.0\times10^{2}$--$8.6\times10^{2}$ at shift
$\epsilon=10^{-2}$ on grids with up to $n=4096$ unknowns, saturating under mesh refinement.  The measured
crossover times $T^{*}=\ln\kappa^{2}/(2\lambda^{+})\approx 7$--$16$ are
short: beyond them the symmetrized recovery wins by 4--26 orders of
magnitude in post-selection cost and, at the longest horizons tested, is
the only working recovery.
% ======================================================================
\subsection{Introduction}
\label{ch3:sec:intro}

Quantum algorithms for partial differential equations (PDEs) divide into
oracle-based constructions, which assume block-encoded or sparse-access
operators \cite{costa2019wave,novikau2023qsvt}, and the explicit-circuit
line initiated for hyperbolic problems by Sato \emph{et
al.}~\cite{sato2024hamiltonian}, in which every gate of the time-stepping
circuit is written down and the cost is dominated by a number of
elementary operations polynomial in the register size.  Explicit-circuit constructions for wave-type equations published to date
operate on \emph{closed} domains: Dirichlet, Neumann, or periodic
boundaries, for which the semi-discrete generator is skew-Hermitian and
the evolution is genuinely
unitary~\cite{sato2024hamiltonian,bosch2025sources}; the extension to
nonconservative systems~\cite{sato2025nonconservative} admits a bulk
damping term $\zeta(x)\,\partial_t u$ via a linear combination of
Hamiltonian simulations, but its boundaries remain closed walls and its
demonstrations are confined to the undamped acoustic and uniform heat
equations.  Scattering and radiation problems,
however, are posed on open domains; truncating the domain with hard walls
reflects every outgoing wave back into the region of interest.  Classical
solvers handle this with absorbing layers, the de facto standard being the
perfectly matched layer (PML) and its convolutional,
complex-frequency-shifted variant
(CPML)~\cite{roden2000cpml,komatitsch2007unsplit,collino2001pml,duru2022pmlreview}.
Absorption is dissipation: the generator acquires a non-trivial Hermitian
part, and the dynamics ceases to be unitary.

Schr\"odingerisation~\cite{jin2023schrodingerisation} is a general route
for running such non-unitary linear dynamics on a quantum computer: the
warped phase-space transform turns $\dot z = Az$ into a genuine
Hamiltonian evolution on one extra register, at the price of a
post-selected recovery slice.  \ref{ch3:fig:pipeline} summarises the
resulting four-step structure of this work.

\begin{figure}[tb]
  \centering
  \begin{tikzpicture}[
      box/.style={draw, rounded corners=2pt, align=center,
                  inner sep=4pt, font=\small, text width=0.86\linewidth},
      lab/.style={font=\scriptsize\itshape, anchor=west},
      node distance=14pt]
    \node[box] (b1) {\textbf{Physics}: second-order acoustic wave
      equation\quad $\partial_t^2 u = c^2\Delta u$};
    \node[box, below=of b1] (b2) {\textbf{First-order reduction}
      (quantum encoding): $\psi=[v;w]$,\;
      $i\,\partial_t\psi = H\psi$ \;---\; unitary};
    \node[box, below=of b2] (b3) {\textbf{Absorbing layer} (CPML,
      convolution memory): $\dot\psi = (-iH-\Lambda)\psi$ \;---\;
      non-Hermitian};
    \node[box, below=of b3] (b4) {\textbf{Schr\"odingerisation}:
      $H_{\mathrm{tot}} = D_\eta\otimes H_1 - I\otimes H_2$ \;---\;
      unitary again ($+\,n_p$ qubits), recover at $p^{*}$};
    \draw[->] (b1) -- node[lab] {\;Ref.~\cite{sato2024hamiltonian}} (b2);
    \draw[->] (b2) -- node[lab] {\;this work, Section~\ref{ch3:sec:generators}} (b3);
    \draw[->] (b3) -- node[lab] {\;Sections~\ref{ch3:sec:schro},
      \ref{ch3:sec:circuits}} (b4);
  \end{tikzpicture}
  \caption{The four-step structure: second-order acoustic physics,
  first-order quantum encoding, convolution absorbing layer applied in
  that first-order form, and Schr\"odingerisation restoring unitarity.
  The cost of step three---the indefiniteness quantified by
  Lemma~\protect\ref{ch3:lem:threshold}---and its remedy (the Lyapunov symmetrizer)
  are the subject of Section~\protect\ref{ch3:sec:symmetrizer}.}
  \label{ch3:fig:pipeline}
\end{figure}  The programme now spans artificial boundary
conditions for the time-dependent Schr\"odinger
equation~\cite{jin2024abc}, inhomogeneous terms and the recovery-threshold
theory~\cite{jin2024recovery}, ill-posed problems~\cite{jin2024illposed},
Maxwell's equations~\cite{jin2023maxwell,ma2024maxwellcircuits}, elastic
waves~\cite{jin2025elastic}, gate-level circuits for parabolic
problems~\cite{hu2024quantum,jin2024heat}, query-optimal
formulations~\cite{jin2025queries}, and implicit-explicit integrators
for multiscale dynamics including the bulk-damped telegraph
equation~\cite{hu2026imex}.  What it does not yet contain is an
absorbing boundary for a classical wave-type equation, at either matrix or
circuit level.

\subsubsection{Related work and positioning}
\label{ch3:sec:related}

The positioning below was verified by a full-text literature review
through June 2026.\footnote{Adversarially cross-checked sweep over
published and preprinted work, 2023--2026; all primary claims below were
confirmed against the full texts of the cited papers.}

\noindent\textbf{The core niche is open.}  No published or preprinted work
through June 2026 combines any absorbing boundary (PML, CPML, complex
absorbing potential, sponge, or Dirichlet-to-Neumann map) for
\emph{classical wave-type equations} with Schr\"odingerisation or any
other unitarisation, at matrix or circuit level.  Jin and
Zhang~\cite{jin2025elastic} (elastic waves) state verbatim that
non-periodic boundaries ``remain an open issue'' left ``to the future.''
Maxwell circuits exist only for perfect-electric-conductor
boundaries~\cite{ma2024maxwellcircuits} and for impedance boundaries at
matrix level~\cite{jin2023maxwell}; the 2026 circuit
papers~\cite{liang2026structure,guseynov2026robin}
(Robin/Neumann/Dirichlet) contain no absorbing-boundary content.
Two 2025 works corroborate this independently.  B\"osch \emph{et
al.}~\cite{bosch2025sources}, simulating lossless wave equations by
Hamiltonian simulation, state that owing to their ``fundamentally
nonenergy conserving nature'' absorbing boundaries and PMLs are ``not
straightforward to implement\ldots within the HS framework,'' and resort
to approximate \emph{random} boundaries that scramble, rather than
absorb, outgoing energy.  Sato \emph{et
al.}~\cite{sato2025nonconservative} admit spatially varying bulk damping
$\zeta(x)\ge0$ in an LCHS formulation, but the resulting Hermitian part
is semidefinite by construction---the diagonal sponge/CAP case that
Lemma~\ref{ch3:lem:threshold} explicitly excludes---no layer-localised
damping or absorbing-layer construction appears, and no damped-wave
circuit is demonstrated.

\noindent\textbf{PML plus Schr\"odingerisation per se is prior art} for the
time-dependent Schr\"odinger equation~\cite{jin2024abc}, at
query-complexity level and without circuits.  Our contribution is
therefore framed as \emph{wave-type equations plus explicit circuits}.

\noindent\textbf{The recovery threshold is established theory}
\cite[Thm.~3.1]{jin2024recovery}, and an indefinite boundary-augmented
Hermitian part was already handled through that threshold for the heat
equation \cite[Remark~1]{jin2024heat}.  Our delta is the
\emph{structural} result: Lemma~\ref{ch3:lem:threshold} proves that the
indefiniteness of CPML memory forms cannot be removed by any diagonal
rescaling and quantifies its rate, $\Theta(\sqrt{\sigma_{\max}c/h})$.

\noindent\textbf{Diagonal similarity preconditioning is prior art.}  Liang and
Liu~\cite{liang2026structure} use a diagonal transform to enforce a
positive-semidefinite Hermitian part before applying a
linear-combination-of-Hamiltonian-simulation algorithm to parabolic
problems.  This is the foil that sharpens our symmetrizer contribution:
Lemma~\ref{ch3:lem:threshold} proves that diagonal transforms are
\emph{insufficient} for CPML generators; the full (non-diagonal) Lyapunov
symmetrizer of Section~\ref{ch3:sec:symmetrizer} is both necessary and
sufficient.

\noindent\textbf{Bell-basis circuits for non-unitary PDEs are prior art}
\cite{hu2024quantum,jin2024heat}.  Our circuit delta is narrower and is
framed accordingly: projector-valued operator strings for separated
multi-field (block) Hamiltonians, the two-level symmetrised splitting, and
the absorbing-boundary application at 17--21 qubits
(Section~\ref{ch3:sec:circuits}).

\subsubsection{Contributions}

Ordered by verified novelty:
\begin{enumerate}
\item the first absorbing boundaries for wave-type equations in any
  unitarisation framework, closing the open problem stated in
  Ref.~\cite{jin2025elastic};
\item the threshold lemma (Lemma~\ref{ch3:lem:threshold}) and the Lyapunov
  symmetrizer that circumvents it, with the diagonal-insufficiency
  contrast to Ref.~\cite{liang2026structure};
\item end-to-end validated explicit circuits, extending the Bell-basis
  constructions of Refs.~\cite{sato2024hamiltonian,hu2024quantum} to
  projector-valued block structure, with gold-standard classical
  validation.
\end{enumerate}

% ======================================================================
\subsection{Background}
\label{ch3:sec:background}

\subsubsection{The wave equation as Hamiltonian simulation}
\label{ch3:sec:wave}

Following Ref.~\cite{sato2024hamiltonian}, the wave equation
$\partial_t^2 u = c^2\Delta u$ is first-order-reduced through
$\psi = [v;\,w]$ with $v=\partial_t u$ and $w = i c\,\partial_x u$
(discretely, $w = icDu$), so that $i\,\partial_t\psi = H\psi$ with
\begin{equation}
  H \;=\; \begin{bmatrix} 0 & D \\ D^{\dagger} & 0 \end{bmatrix},
  \label{ch3:eq:waveH}
\end{equation}
where $D$ is the forward-difference operator and
$D^{\dagger}=-D_{\mathrm{b}}$ its negative backward counterpart.  Mixed
Dirichlet/Neumann boundary conditions emerge from the operator products
(the conventions of Ref.~\cite{sato2024hamiltonian}, verified numerically
in our implementation).  In two dimensions we use the separated-block form
$\psi = [v;\,w_x;\,w_y]$---one block per field, in the style of Eq.~(32)
of Ref.~\cite{sato2024hamiltonian} with $d=3$---instead of the
complex-combined two-block form, because per-direction damping requires
separated gradient components.

\subsubsection{The convolutional perfectly matched layer}
\label{ch3:sec:cpml}

The CPML stretches each coordinate by the complex-frequency-shifted (CFS)
factor
\begin{equation}
  s(\omega) \;=\; \kappa + \frac{\sigma}{\alpha + i\omega},
  \label{ch3:eq:cfs}
\end{equation}
which in the time domain replaces every spatial derivative by a causal
convolution with an exponential kernel~\cite{roden2000cpml},
\begin{equation}
  \partial_x f \;\longrightarrow\; \frac{1}{\kappa}\,\partial_x f
    + \zeta * \partial_x f, \qquad
  \zeta(t) = -\frac{\sigma}{\kappa^{2}}\,
    e^{-(\sigma/\kappa+\alpha)t}\,\theta(t).
  \label{ch3:eq:kernel}
\end{equation}
Each convolution is carried by one memory field obeying the local ODE
\begin{equation}
  \partial_t \phi \;=\; -\Bigl(\frac{\sigma}{\kappa}+\alpha\Bigr)\phi
    \;-\; \frac{\sigma}{\kappa^{2}}\,\partial_x f .
  \label{ch3:eq:memoryode}
\end{equation}
The damping profile is polynomially graded inside a layer of
$n_{\mathrm{pml}}$ points, $\sigma(d) = \sigma_{\max}(d/L)^{m}$ with
$\sigma_{\max} = -(m{+}1)\,c\ln R_0/(2L)$~\cite{collino2001pml}; we use
the standard recipe $m=2$, $R_0=10^{-3}$ with the parameter conventions of
Komatitsch and Martin~\cite{komatitsch2007unsplit}.  The profile is
sampled at the staggered locations where each field effectively lives
($v$ at integer nodes, $w$ at half-cells).  This stagger is load-bearing:
an ablation shows that unstaggered sampling degrades the measured
reflection by factors of 120--235.

\subsubsection{Schr\"odingerisation}
\label{ch3:sec:schro}

Write the semi-discrete absorbing dynamics as $\dot z = Az$ and split
$A = H_1 + iH_2$ with $H_1=(A+A^{\dagger})/2$ and
$H_2=(A-A^{\dagger})/(2i)$, both Hermitian.  Only $H_1$ obstructs
unitarity: $\tfrac{d}{dt}\|z\|^{2} = 2\,z^{\dagger}H_1z$, so $e^{At}$ is
unitary precisely when $H_1=0$, whereas $iH_2$ is already skew-Hermitian.
The warped transform $w(t,p)=e^{-p}z(t)$ trades that obstruction for one
extra dimension: because $\partial_p w = -w$, the offending term can be
rewritten as $H_1 w = -H_1\partial_p w$, giving
\begin{equation}
  \partial_t w \;=\; -H_1\partial_p w + iH_2 w .
  \label{ch3:eq:warped}
\end{equation}
The rewrite is applied to the $H_1$ term alone---$iH_2$ needs no
repair---and that is precisely what it buys: $\partial_p$ is
skew-Hermitian, hence so is $-H_1\otimes\partial_p$, and the enlarged
generator is skew-Hermitian in full.  A Fourier series in $p$, under
which $\partial_p$ acts as multiplication by $i\eta_k$ with $\eta_k$ the
$k$th grid frequency, decouples Eq.~\eqref{ch3:eq:warped} into unitary
evolutions per mode~\cite{jin2023schrodingerisation},
\begin{equation}
  \hat w_k(T) \;=\; e^{-iT(\eta_k H_1 - H_2)}\,\hat w_k(0),
  \label{ch3:eq:permode}
\end{equation}
i.e.\ a Hamiltonian simulation of
$H_{\mathrm{tot}} = D_\eta\otimes H_1 - I\otimes H_2$ on the system
register plus one $p$ register of $n_p$ qubits.  Since $e^{-p}$ is not
square-integrable on $p<0$, the profile is extended to the left half-line
by a normalisable $g$; the initial data $w(0,p)=g(p)\,z(0)$ is then a
product state across the two registers.  The
solution is recovered from a slice,
\begin{equation}
  z(T) \;=\; e^{p^{*}} w(T,p^{*}),
  \qquad
  p^{*} \;\ge\; \lambda^{+}_{\max}(H_1)\,T,
  \label{ch3:eq:recovery}
\end{equation}
where the threshold on $p^{*}$ is Theorem~3.1 of
Ref.~\cite{jin2024recovery}.  The kinked profile $g(p)=e^{-|p|}$ is only
$O(\Delta p)$ accurate; the $C^{1}$ cubic profile of
\cite[Remark~4.9]{jin2024recovery} restores $O(\Delta p^{2})$ while
keeping $g(p)=e^{-p}$ for $p\ge0$, leaving Eq.~\eqref{ch3:eq:recovery}
unchanged.

% ======================================================================
\subsection{Absorbing generators}
\label{ch3:sec:generators}

\subsubsection{One dimension: exact collapse of the convolution}
\label{ch3:sec:1d}

With $\kappa=1$, $\alpha=0$ the convolution collapses \emph{exactly} (for
data supported outside the layer) to block-diagonal damping,
\begin{equation}
  \dot\psi = A\psi, \qquad
  A \;=\; -iH - \diag(\sigma_v,\, \sigma_w),
  \label{ch3:eq:1dcollapse}
\end{equation}
so $H_1 = -\diag(\sigma)\preceq 0$ and the recovery
threshold~\eqref{ch3:eq:recovery} is unconditional ($\lambda^{+}=0$, any small
$p^{*}>0$ is valid).  The proof is the memory ansatz
$\phi = -\sigma v/\gamma$, which satisfies the memory
ODE~\eqref{ch3:eq:memoryode} identically for data supported outside the
layer; we verify the collapse numerically to $6\times10^{-13}$.  The
general CFS memory form (two extra block qubits) is provided for
$\kappa>1$ or $\alpha>0$.

\subsubsection{Two dimensions: memory form and classical validation}
\label{ch3:sec:2d}

In 2D the damping acts per direction, so the state uses the
separated-block form padded to eight blocks (three extra block qubits),
\begin{equation}
  \bigl[\,v;\; w_x;\; w_y;\; \phi_{vx};\; \phi_{vy};\;
        \phi_{x};\; \phi_{y};\; 0\,\bigr],
  \label{ch3:eq:blocks}
\end{equation}
with the memory injections masked to the layer strips.  Corners need no
special treatment---both directions' memory fields are simultaneously
active, the structural advantage of unsplit CPML over split
formulations~\cite{roden2000cpml}.  Long-time stability is a serious
concern for PMLs in
general~\cite{becache2003stability,becache2004longtime,duru2022pmlreview};
here we measure $\max\operatorname{Re}\operatorname{eig}(A) = 0$ to
round-off across all tested configurations.

\emph{Gold-standard validation.}  Comparing the truncated domain against
a $4\times$ enlarged reference, the true reflection is $3.4\times10^{-4}$
(1D, $n_{\mathrm{pml}}=12$) and $2.8\times10^{-4}$ (2D $32\times32$,
$n_{\mathrm{pml}}=8$), at the floor reported by Komatitsch and
Martin~\cite{komatitsch2007unsplit}.

\emph{Convergence} [1D gold standard, $N=128$ versus an $N=512$
reference, plateau over $T\in\{100,120,140\}$ with plateau spread
$<10^{-9}$, i.e.\ fully converged; \ref{ch3:fig:convergence}(a,b)]:
\begin{itemize}
\item versus $n_{\mathrm{pml}}$ at $R_0=10^{-3}$: $6.7\times10^{-4}$ (4),
  $2.6\times10^{-4}$ (6), $2.6\times10^{-4}$ (8), $3.4\times10^{-4}$ (12),
  $4.2\times10^{-4}$ (16)---\emph{non-monotonic}, with a minimum at
  $n_{\mathrm{pml}}\approx6$--$8$.  The two contributions are the
  discrete transition error, which decreases with layer width, and the
  wall-return reflection, which \emph{increases toward the design value
  from below}: node sampling of the polynomial profile overestimates
  $\int\sigma\,dx$ by a factor $1+O(1/n_{\mathrm{pml}})$, so the
  effective round-trip reflection is
  $R_{\mathrm{eff}}=R_0^{\,\Sigma\sigma h/\!\int\!\sigma}$
  ($R_0^{1.38}=7.1\times10^{-5}$ at $n_{\mathrm{pml}}=4$ rising to
  $R_0^{1.09}=5.2\times10^{-4}$ at $16$); for
  $n_{\mathrm{pml}}\ge8$ the measured reflection equals this
  prediction within $5$--$20\%$, i.e.\ the layer performs at design,
  converging to $R_0$ from below [\ref{ch3:fig:convergence}(a)].
  Jointly rescaling the staggered profile pair so the discrete sum
  matches the design integral (\texttt{calibration='discrete'},
  Appendix~\ref{ch3:appD:validation}) makes the realised reflection flat at
  $0.81$--$1.08\,R_0$ over the whole sweep---a one-line calibration fix
  that applies to classical CPML codes as well;
\item versus $R_0$ at $n_{\mathrm{pml}}=12$: $4.5\times10^{-3}$
  ($10^{-2}$), $3.4\times10^{-4}$ ($10^{-3}$), $4.3\times10^{-5}$
  ($10^{-4}$)---tracking the design reflection $R_0$ within a factor
  ${\sim}2.3$ over two decades.  The Collino--Tsogka amplitude
  formula~\cite{collino2001pml} is predictive in this discretisation.
\end{itemize}

\begin{figure}
  \includegraphics[width=\linewidth]{chapter3/paper_P1_convergence.png}
  \caption{Classical and Schr\"odingerisation convergence.
  (a)~Gold-standard reflection (truncated domain versus $4\times$
  reference) versus layer width $n_{\mathrm{pml}}$ at design reflection
  $R_0=10^{-3}$.  The curve is non-monotonic: the transition error falls
  with layer width while the wall return rises toward the design $R_0$
  from below, because node sampling of the polynomial profile
  over-damps thin layers (dashed curve: the predicted discrete
  wall-return $R_{\mathrm{eff}}=R_0^{\,\Sigma\sigma
  h/\!\int\!\sigma}$, which the measurement tracks within
  $5$--$20\%$ for $n_{\mathrm{pml}}\ge8$).  Red squares: the
  discretely calibrated profile pair (Appendix~\protect\ref{ch3:appD:validation}),
  for which the realised reflection is flat at the design value.
  (b)~Reflection versus $R_0$ at $n_{\mathrm{pml}}=12$: the measured
  reflection tracks the Collino--Tsogka design value within a factor
  ${\sim}2.3$ over two decades.
  (c)~Recovery error versus $p$-register size $n_p$ for the kinked
  $e^{-|p|}$ warping profile [$O(\Delta p)$] and the $C^{1}$ cubic
  profile [$O(\Delta p^{2})$] of Ref.~\protect\cite{jin2024recovery}, 1D CPML
  benchmark.
  (d)~Trotter splitting error versus step count at orders 1 and 2
  (fitted slopes $-1.02$ and $-2.01$), measured against the exact
  Schr\"odingerised evolution to isolate the splitting error.}
  \label{ch3:fig:convergence}
\end{figure}

\subsubsection{The recovery-threshold obstruction}
\label{ch3:sec:lemma}

The memory form pays a price.  The injection column ($\gamma$ on one
side) and the memory-source row ($\sigma\partial_x/\gamma$ on the other)
of each memory field are structurally asymmetric, and this asymmetry
cannot be scaled away.  Throughout this section
$\beta_j = \sigma_j/\kappa_j+\alpha_j\ge0$ denotes the CFS decay rate at
node $j$; we set the wave speed $c=1$ (so the stencil entries obey
$|D_{ji}|=1/h$), the general rate restoring a factor of $c$.

% Reconciled 2026-06-11: the (w_i,phi_j) coupling is one-directional
% (A[w,phi]=0), so the entry s lands in H_1 as s/2 and the eigenvalue of
% [[0,s/2],[s/2,-beta]] is (sqrt(beta^2+s^2)-beta)/2 — both pairs are
% governed by the same f.  The lemma statement now matches the proof.
\begin{lemma}[diagonal irreducibility]
\label{ch3:lem:threshold}
Let $A$ be the memory-form CPML generator of Section~\ref{ch3:sec:2d} and let
$G=\diag(\gamma_j)>0$ be any diagonal rescaling of the memory fields.
Let $s_j$ denote the magnitude of the memory-source coupling
($s_j\sim\sigma_j/(\gamma_j h)$ for the difference stencil).  Then the
Hermitian part $H_1(G)$ of the rescaled generator satisfies
\begin{equation}
  \lamax\bigl(H_1(G)\bigr) \;\ge\;
  \max_j\, \max\bigl\{\, f(\gamma_j,\beta_j),\; g(s_j,\beta_j) \,\bigr\}
  \;>\; 0,
  \label{ch3:eq:lemmabound}
\end{equation}
with $f(\gamma,\beta) = \bigl(\sqrt{\beta^{2}+\gamma^{2}}-\beta\bigr)/2$,
arising from the $2\times2$ principal submatrix
$\bigl[\begin{smallmatrix} 0 & \gamma/2 \\ \gamma/2 & -\beta
\end{smallmatrix}\bigr]$
on a $(v_j,\phi_j)$ pair by Cauchy interlacing, and
$g(s,\beta) = f(s,\beta)$ from a $(w_i,\phi_j)$ pair (the one-directional
coupling entry $s$ enters $H_1$ as $s/2$, so both pairs are governed by
the same $f$).  Optimising $\gamma$ against the two bounds gives
\begin{equation}
  \lamax(H_1) \;=\; \Theta\!\left(\sqrt{\sigma_{\max}\,c/h}\right)
  \qquad \text{as } h\to0,
  \label{ch3:eq:lemmarate}
\end{equation}
independent of the scaling.
\end{lemma}

The full proof is given in Appendix~\ref{ch3:app:proof}.  Its consistency
check is instructive: the collapsed 1D form of Section~\ref{ch3:sec:1d} has
\emph{no} injection pair (the memory is eliminated and the damping is
diagonal), so the lemma does not apply---and indeed
$H_1=-\diag(\sigma)\preceq0$ there.

\emph{Consequence.}  By the threshold of Eq.~\eqref{ch3:eq:recovery}
\cite[Thm.~3.1]{jin2024recovery}, recovery requires
$p^{*}\ge\lambda^{+}T$ with $\lambda^{+}\equiv\lamax^{+}(H_1)$, and the
post-selection probability of the recovery window scales as
$e^{-2p^{*}} = e^{-2\lambda^{+}T}$.  Numerically: at $\sigma_{\max}=0.5$
on an $8\times8$ grid, $\lambda^{+}=0.683$; recovery below the threshold
is $O(1)$ wrong, while above it the error drops to $2\times10^{-3}$
relative ($10^{-4}$ with $n_p=11$); see \ref{ch3:fig:threshold}.  The
sponge/CAP alternative (diagonal damping $\sigma_x+\sigma_y$, four
blocks) has $H_1\preceq0$ unconditionally at the price of matched-layer
quality.

% ----------------------------------------------------------------------
% CPML versus sponge at matched resources (end of Sec. III "The
% recovery-threshold obstruction"; replaces the [[CPML vs sponge at
% matched resources --- pending sweeps]] todo).
% Source of every number: paper/experiments/cpml_vs_sponge.py
%   -> paper/experiments/cpml_vs_sponge_results.json (run 2026-06-11).
% Labels prefixed cmp: to avoid collisions.
% ----------------------------------------------------------------------

\emph{CPML versus sponge at matched resources.}  A matched-resource
sweep quantifies this tradeoff: both absorbers on the same $32\times32$
grid with equal layer widths $n_{\mathrm{pml}}\in\{4,8,12\}$ and
default parameters, each measured against one shared $4\times$ enlarged
($128\times128$) hard-wall reference by the gold-standard protocol of
Section~\ref{ch3:sec:2d}, the reflection being the maximum interior error over
$T\in\{8,16,24,32,40\}$ (\ref{ch3:cmp:tab:matched}).  At equal layer
budget the CPML reaches its discrete wall-return level of
$2.7\times10^{-4}$ by $n_{\mathrm{pml}}=8$ [consistent with the
$2.8\times10^{-4}$ gold-standard value of Section~\ref{ch3:sec:2d}], while the
sponge stays at the percent level, improving only from
$6.2\times10^{-2}$ to $3.6\times10^{-2}$ across the sweep---a quality
gap of a factor of $170$ at $n_{\mathrm{pml}}=8$ ($130$ at $12$).  The
quantum-side cost inverts the ranking.  The sponge has $\lambda^{+}=0$
exactly, so its recovery post-selects at the $T$-independent
$e^{2p^{*}}\approx1.2$ (default slice $p^{*}\approx0.11$); the CPML
pays $e^{2\lambda^{+}T}$, with $\lambda^{+}=0.93$--$1.54$ tracking the
$\Theta(\sqrt{\sigma_{\max}c/h})$ rate of Lemma~\ref{ch3:lem:threshold} to
within $5\%$---already $7\times10^{7}$ at $T=8$ and $1.6\times10^{39}$
at $T=40$ for $n_{\mathrm{pml}}=8$.  Neither absorber dominates at
matched resources: the sponge trades two orders of magnitude in
reflection for an essentially free recovery, and the CPML trades an
exponentially growing post-selection penalty for matched-layer
accuracy---precisely the penalty that the symmetrizer of
Section~\ref{ch3:sec:symmetrizer} replaces by a $T$-independent conditioning
factor.

\begin{table}[tb]
\caption{CPML versus sponge at matched resources: $32\times32$ grid,
equal layer width $n_{\mathrm{pml}}$, default parameters [$m=2$,
$R_0=10^{-3}$, Collino--Tsogka amplitudes
$\sigma_{\max}=2.59/1.30/0.86$ at $n_{\mathrm{pml}}=4/8/12$].
Reflection is the gold-standard interior error against a single shared
$4\times$ enlarged ($128\times128$) hard-wall reference, maximised over
$T\in\{8,16,24,32,40\}$ ($\lVert z_0\rVert=1$).
$\lambda^{+}=\lamax^{+}(H_1)$ sets the CPML post-selection penalty
$e^{2\lambda^{+}T}$ (quoted at $T=40$) and tracks
$\sqrt{\sigma_{\max}}$ to within $5\%$; the sponge has
$\lambda^{+}=0$, post-selecting at the $T$- and
$n_{\mathrm{pml}}$-independent $e^{2p^{*}}=1.2$ for the default
recovery slice $p^{*}=0.11$.}
\label{ch3:cmp:tab:matched}
\begin{tabular}{lllll}
\toprule
 & \multicolumn{2}{c}{Reflection (max over $T$)} &
   \multicolumn{2}{c}{CPML recovery cost} \\
\cline{2-3}\cline{4-5}
$n_{\mathrm{pml}}$ & CPML & sponge & $\lambda^{+}(H_1)$ &
$e^{2\lambda^{+}T}$, $T{=}40$ \\
\midrule
4  & $4.2\times10^{-3}$ & $6.2\times10^{-2}$ & 1.54 & $2.2\times10^{53}$ \\
8  & $2.7\times10^{-4}$ & $4.7\times10^{-2}$ & 1.13 & $1.6\times10^{39}$ \\
12 & $2.8\times10^{-4}$ & $3.6\times10^{-2}$ & 0.93 & $1.5\times10^{32}$ \\
\bottomrule
\end{tabular}
\end{table}


\begin{figure}
  \includegraphics[width=\linewidth]{chapter3/fig4_2d_recovery_threshold.png}
  \caption{The recovery threshold in action (2D CPML, $8\times8$ grid,
  $\sigma_{\max}=0.5$, $T=10$).  Relative recovery error versus the slice
  position $p^{*}$: below the threshold $p^{*}=\lambda^{+}T$ (dashed
  line, $\lambda^{+}=0.683$) the recovered state is $O(1)$ wrong; just
  above it the error drops by four orders of magnitude, then grows
  slowly as $e^{p^{*}}$ amplifies the $p$-grid discretisation floor.}
  \label{ch3:fig:threshold}
\end{figure}

\subsubsection{The Lyapunov symmetrizer}
\label{ch3:sec:symmetrizer}

A matrix $W>0$ solving the Lyapunov inequality
$WA + A^{\dagger}W\preceq0$ gives, with $S=W^{1/2}$, a transformed
generator $\At = SAS^{-1}$ whose Hermitian part satisfies
\begin{equation}
  \Ht \;=\; \tfrac12\, S^{-1}\bigl(WA + A^{\dagger}W\bigr)S^{-1}
  \;\preceq\; 0
  \label{ch3:eq:congruence}
\end{equation}
by congruence.  Recovery of $\tilde z = Sz$ is then valid at any
$p^{*}>0$, and $z = S^{-1}\tilde z$ is recovered classically (or folded
into the readout).  Since $A$ is only marginally stable (interior modes
are neutral), we construct $W$ from the $\epsilon$-shifted generator:
\begin{equation}
  A_\epsilon = A - \epsilon I, \qquad
  A_\epsilon^{\dagger}W + WA_\epsilon = -I, \qquad
  S = W^{1/2},
  \label{ch3:eq:lyapunov}
\end{equation}
the shift being undone classically by the factor $e^{\epsilon T}$.  An
exact identity follows (Proposition~\ref{ch3:appE:prop} in
Appendix~\ref{ch3:appE:symmetrizer}: congruence gives
$\mathrm{Herm}(SA_\epsilon S^{-1})=-W^{-1}/2$ exactly; verified to better
than $10^{-11}$ relative in all 18 measured configurations): the
transform of the \emph{original} $A$ has
\begin{equation}
  \lamax(\Ht) \;=\; \epsilon - \frac{1}{2\lamax(W)} \;\approx\; \epsilon,
  \label{ch3:eq:identity}
\end{equation}
while the transform of $A_\epsilon$ has $\Ht\preceq0$ strictly.  We
therefore use the shifted transform; the undo factor
$e^{\epsilon T}\le e^{0.4}$ over the horizons tested is trivial.

\begin{table}
\caption{Lyapunov-symmetrizer conditioning across all measured
configurations [grids $(n_x,n_y)$ qubits per dimension, i.e.\ $8\times8$,
$16\times8$, $16\times16$ points; $n$ is the full state dimension
including the eight blocks].  $\lambda^{+}=\lamax^{+}(H_1)$ is the
indefiniteness of the untransformed generator;
$\kappa_2(S)=\sqrt{\operatorname{cond}_2(W)}$ the conditioning of the
symmetrizer; $\lamax(\Ht^{(\epsilon)})$ the (strictly negative) largest
Hermitian-part eigenvalue of the $\epsilon$-shifted transform
$SA_\epsilon S^{-1}$; and $T^{*}=\ln\kappa^{2}/(2\lambda^{+})$ the
crossover time beyond which the symmetrizer beats the direct
above-threshold recovery.  All Lyapunov residuals are
$\le1.4\times10^{-14}$ (relative).}
\label{ch3:tab:kappa}
\begin{tabular}{llllllll}
\toprule
Grid & $n$ & $\sigma_{\max}$ & $\epsilon$ & $\lambda^{+}(H_1)$ &
$\kappa_2(S)$ & $\lamax(\Ht^{(\epsilon)})$ & $T^{*}$ \\
\midrule
$8\times8$ & 512 & 0.5 & $10^{-2}$ & 0.683 & $4.0\times10^{2}$ & $-4.6\times10^{-6}$ & 8.8 \\
$8\times8$ & 512 & 0.5 & $10^{-3}$ & 0.683 & $3.0\times10^{3}$ & $-8.1\times10^{-8}$ & 11.7 \\
$8\times8$ & 512 & 0.5 & $10^{-4}$ & 0.683 & $1.1\times10^{4}$ & $-6.0\times10^{-9}$ & 13.6 \\
$8\times8$ & 512 & 1.0 & $10^{-2}$ & 0.927 & $5.1\times10^{2}$ & $-5.2\times10^{-6}$ & 6.7 \\
$8\times8$ & 512 & 1.0 & $10^{-3}$ & 0.927 & $2.9\times10^{3}$ & $-1.6\times10^{-7}$ & 8.6 \\
$8\times8$ & 512 & 1.0 & $10^{-4}$ & 0.927 & $9.8\times10^{3}$ & $-1.4\times10^{-8}$ & 9.9 \\
$16\times8$ & 1024 & 0.5 & $10^{-2}$ & 0.687 & $5.1\times10^{2}$ & $-3.0\times10^{-6}$ & 9.1 \\
$16\times8$ & 1024 & 0.5 & $10^{-3}$ & 0.687 & $7.0\times10^{3}$ & $-1.5\times10^{-8}$ & 12.9 \\
$16\times8$ & 1024 & 0.5 & $10^{-4}$ & 0.687 & $3.1\times10^{4}$ & $-8.0\times10^{-10}$ & 15.0 \\
$16\times8$ & 1024 & 1.0 & $10^{-2}$ & 0.956 & $7.5\times10^{2}$ & $-2.4\times10^{-6}$ & 6.9 \\
$16\times8$ & 1024 & 1.0 & $10^{-3}$ & 0.956 & $7.0\times10^{3}$ & $-2.8\times10^{-8}$ & 9.3 \\
$16\times8$ & 1024 & 1.0 & $10^{-4}$ & 0.956 & $2.7\times10^{4}$ & $-1.8\times10^{-9}$ & 10.7 \\
$16\times16$ & 2048 & 0.5 & $10^{-2}$ & 0.694 & $5.7\times10^{2}$ & $-2.4\times10^{-6}$ & 9.1 \\
$16\times16$ & 2048 & 0.5 & $10^{-3}$ & 0.694 & $9.9\times10^{3}$ & $-7.7\times10^{-9}$ & 13.3 \\
$16\times16$ & 2048 & 0.5 & $10^{-4}$ & 0.694 & $6.6\times10^{4}$ & $-1.7\times10^{-10}$ & 16.0 \\
$16\times16$ & 2048 & 1.0 & $10^{-2}$ & 0.972 & $8.6\times10^{2}$ & $-1.9\times10^{-6}$ & 6.9 \\
$16\times16$ & 2048 & 1.0 & $10^{-3}$ & 0.972 & $1.2\times10^{4}$ & $-8.7\times10^{-9}$ & 9.7 \\
$16\times16$ & 2048 & 1.0 & $10^{-4}$ & 0.972 & $6.3\times10^{4}$ & $-3.4\times10^{-10}$ & 11.4 \\
\bottomrule
\end{tabular}
\end{table}

\begin{table}
\caption{End-to-end recovery tradeoff: direct versus symmetrized
recovery, errors
$\lVert z_{\mathrm{rec}}-e^{AT}z_0\rVert_2/\lVert e^{AT}z_0\rVert_2$.
``Small $p^{*}$'' is the library-default slice three grid points above
zero.  For the untransformed system above threshold,
$p^{*}=\lambda^{+}T+1$ and the post-selection penalty is $e^{2p^{*}}$;
entries marked ``--'' exceed the $p$-domain
($\lambda^{+}T+1>p_{\max}$).  The symmetrized column uses the
$\epsilon$-shifted transform at $\epsilon=10^{-3}$ recovered at the same
small $p^{*}$, with the $T$-independent penalty
$\kappa_2^{2}(S)$ of \protect\ref{ch3:tab:kappa}.  The $16\times16$ run uses a
reduced $p$ register ($n_p=8$), which sets its higher error floor.}
\label{ch3:tab:tradeoff}
\resizebox{\linewidth}{!}{%
\begin{tabular}{llllllll}
\toprule
 & & \multicolumn{4}{c}{Untransformed} & \multicolumn{1}{c}{Symmetrized ($\epsilon=10^{-3}$)} \\
\cline{3-6}\cline{7-7}
Configuration & $T$ & err.\ (small $p^{*}$) & $p^{*}\!\ge\!\lambda^{+}T$ &
err.\ (above thr.) & penalty $e^{2p^{*}}$ & err.\ (small $p^{*}$) \\
\midrule
$8\times8$ ($\sigma_{\max}=0.5$, $n_p=11$) & 10 & $3.4$ & 7.8 & $1.2\times10^{-4}$ & $6.3\times10^{6}$ & $1.2\times10^{-6}$ \\
 & 20 & $8.3$ & 14.7 & $2.4\times10^{-1}$ & $5.4\times10^{12}$ & $2.4\times10^{-6}$ \\
 & 40 & $2.6\times10^{1}$ & -- & -- & -- & $4.4\times10^{-6}$ \\
$8\times8$ ($\sigma_{\max}=1.0$, $n_p=11$) & 10 & $7.2$ & 10.3 & $2.2\times10^{-2}$ & $8.4\times10^{8}$ & $2.1\times10^{-5}$ \\
 & 20 & $1.6\times10^{1}$ & 19.6 & $3.0\times10^{2}$ & $1.0\times10^{17}$ & $2.3\times10^{-5}$ \\
 & 40 & $2.9\times10^{1}$ & 38.1 & $7.3\times10^{15}$ & $1.2\times10^{33}$ & $3.5\times10^{-5}$ \\
$16\times8$ ($\sigma_{\max}=0.5$, $n_p=10$) & 10 & $2.5$ & 7.9 & $1.1\times10^{-4}$ & $7.4\times10^{6}$ & $4.2\times10^{-6}$ \\
 & 20 & $4.8$ & 14.8 & $4.5\times10^{-1}$ & $6.7\times10^{12}$ & $6.0\times10^{-6}$ \\
 & 40 & $1.3\times10^{1}$ & -- & -- & -- & $1.1\times10^{-5}$ \\
$16\times16$ ($\sigma_{\max}=0.5$, $n_p=8$) & 10 & $2.0$ & 8.0 & $1.3\times10^{-2}$ & $9.2\times10^{6}$ & $1.1\times10^{-4}$ \\
 & 20 & $5.6$ & 15.0 & $7.5\times10^{1}$ & $1.0\times10^{13}$ & $5.2\times10^{-4}$ \\
 & 40 & $2.5\times10^{1}$ & -- & -- & -- & $3.8\times10^{-3}$ \\
\bottomrule
\end{tabular}%
}
\end{table}

\emph{Measured results} [grids $8\times8$, $16\times8$, $16\times16$
($n=512$--$2048$), $\sigma_{\max}\in\{0.5,1.0\}$,
$\epsilon\in\{10^{-2},10^{-3},10^{-4}\}$; Lyapunov residuals
$\le1.4\times10^{-14}$; \ref{ch3:tab:kappa} and \ref{ch3:tab:tradeoff},
\ref{ch3:fig:symmetrizer}]:
\begin{itemize}
\item \emph{End-to-end recovery at the default small $p^{*}$}
  ($\epsilon=10^{-3}$): transformed errors
  $1.2\times10^{-6}$--$3.5\times10^{-5}$ on the configurations with
  $n_p\in\{10,11\}$, essentially flat in $T$ over $T\in\{10,20,40\}$; the
  $16\times16$ run with a reduced $p$ register ($n_p=8$) gives
  $1.1\times10^{-4}$--$3.8\times10^{-3}$, consistent with its
  $O(\Delta p^{2})$ floor.  The untransformed system at the same $p^{*}$
  returns $O(1)$--$O(10)$ garbage growing with $T$.
\item \emph{The direct above-threshold method is worse than its nominal
  $e^{2\lambda^{+}T}$ cost}: at $T=20$ it degrades (errors $0.24$--$0.45$
  at $\sigma_{\max}=0.5$) and at $T=40$ it explodes (error up to
  $7\times10^{15}$ at $p^{*}=38$) because $e^{p^{*}}$ amplifies the
  $p$-grid floor.  At long $T$ the symmetrizer is not merely cheaper---it
  is the only working recovery.
\item \emph{Conditioning}:
  $\kappa_2(S)=4.0\times10^{2}$--$8.6\times10^{2}$ at $\epsilon=10^{-2}$
  (the recommended setting), flat in $\sigma_{\max}$, scaling as
  ${\sim}\epsilon^{-0.6\,\text{to}\,-1.0}$---and \emph{saturating} in
  $n$: extending the measurement to $n=4096$ gives
  $\kappa_2(S)=564.5$ at $\epsilon=10^{-2}$ versus $566.9$ at $n=2048$
  (local slope $-0.006$; over the full range $n=512$--$4096$ the fitted
  power-law exponent drops from $0.25$ to $0.16$ with degrading fit
  quality, i.e.\ a power law is the wrong model because $\kappa$ is
  flattening).  $\lambda^{+}$ itself saturates ($0.683\to0.694$ over
  $8\times$ in $n$) and $\lambda_{\min}(W)$ stays at $0.66$--$0.68$, so
  the $\kappa$ growth is driven entirely by
  $\lambda_{\max}(W)\sim1/(2\epsilon)$: the conditioning is
  $\epsilon$-controlled, not grid-controlled, and the symmetrizer cost
  does not blow up under mesh refinement.  The eigenvector symmetrizer
  $W=(VV^{\dagger})^{-1}$ is four to six orders of magnitude worse
  conditioned ($\kappa_2\sim3\times10^{8}$) and fails outright on
  near-defective configurations---the Lyapunov route is the right
  construction.
\item \emph{Tradeoff}:
  $T^{*}=\ln\kappa^{2}/(2\lambda^{+})=6.7$--$16.0$ across all
  configurations; at $T=20$ the symmetrizer wins by 4--9 orders of
  magnitude in post-selection cost, at $T=40$ by 16--26 orders.  $T^{*}$
  grows only logarithmically with $\kappa$, hence very weakly with grid
  size.
\end{itemize}

\emph{Structure and compressibility of $S$.}  We measured the structure
of $S$ directly and tested 25 approximate symmetrizers (distance-banded,
magnitude-thresholded, block-pattern, identity-plus-low-rank, hybrid,
and inverse-side-banded families) for dissipativity and end-to-end
recovery on the $n=512$ and $n=1024$ systems.  Three findings.
(i)~Identity-plus-low-rank is \emph{ruled out}: the correction
$E=S-\alpha I$ is numerically full rank (995 of 1024 singular values
above $10^{-2}$ of the top), and $\lamax(\tilde H_1)$ stays at the
\emph{untransformed} $\lambda^{+}=0.67$ for ranks $8$--$64$---no
few-term LCU implementation of $S$ exists.  (ii)~Strict sparsification
preserving $\lamax(\tilde H_1)\le\epsilon$ requires essentially the
full active support ($\approx48$--$51\%$ of entries; on the remaining
``dead'' indices---interior memory and padding, $30$--$50\%$ of the
space---$S$ is exactly a multiple of the identity, which is free
structure).  (iii)~\emph{Moderate compression works} once the residual
$\lambda_b$ of an approximate $S_b$ is absorbed as an additional shift
(undone by $e^{\lambda_b T}\le e^{2.6}\approx14$ at the worst measured
$\lambda_b=0.066$, $T=40$): magnitude thresholding at $35\%$ density
recovers to $8.7\times10^{-6}$ ($T=10$) and $5.4\times10^{-5}$
($T=40$); a rank-$32$-plus-band-$8$ hybrid at $26\%$ reaches
$1.3\times10^{-5}$ and $6.6\times10^{-4}$ (all $n=1024$; the
conditioning of every usable $S_b$ stays ${\sim}5\times10^{2}$, equal
to the exact $S$).  The energy of $E$ is $99.6\%$ concentrated on
layer-strip index pairs.  The implementation route is therefore
block-encoding of a ${\sim}35\%$-dense operator (times the free
dead-index identity factor), not a local circuit or a short LCU; at
current theory the symmetrizer is best used as classical pre/\allowbreak
post-processing, and a structured-ansatz Lyapunov construction (solving
for $W$ within a polynomially parametrised operator family) is the open
problem these measurements sharpen.

\emph{Honest limitations.}  The construction requires a dense $O(n^{3})$
classical precomputation of $W$, $S$, $S^{-1}$; the fate of the sparse
quantised-tensor-train (QTT) term structure under the dense similarity
(Trotter circuits for $\At$) is not assessed; the $\kappa$ scaling
beyond $n=4096$ is extrapolation (though saturating); the compression
study is at $n\le1024$.

\begin{figure}
  \includegraphics[width=\linewidth]{chapter3/paper_P2_symmetrizer.png}
  \caption{The Lyapunov symmetrizer.  (a)~Relative recovery error versus
  evolution time $T$ for the 2D CPML system ($8\times8$ grid,
  $\sigma_{\max}=0.5$): untransformed recovery at the default small
  $p^{*}$ (squares) is $O(1)$ garbage; untransformed recovery above the
  threshold $p^{*}\ge\lambda^{+}T$ (diamonds) degrades with $T$ as
  $e^{p^{*}}$ amplifies the $p$-grid floor and leaves the admissible
  $p$-domain entirely by $T=40$; the Lyapunov-symmetrized recovery at the
  same small $p^{*}$ (circles, $\epsilon=10^{-3}$, shifted transform)
  stays at the $10^{-6}$ level, essentially flat in $T$.
  (b)~Post-selection cost on the $16\times16$ grid: the direct
  $e^{2\lambda^{+}T}$ penalty (solid) against the $T$-independent
  symmetrizer penalty $\kappa_2^{2}(S)$ at $\epsilon=10^{-2}$ (dashed)
  and $10^{-3}$ (dotted); the crossovers are the $T^{*}$ values of
  \protect\ref{ch3:tab:kappa}.}
  \label{ch3:fig:symmetrizer}
\end{figure}

% ======================================================================
\subsection{Explicit circuits}
\label{ch3:sec:circuits}

\subsubsection{Projector-extended Bell-basis term evolution}
\label{ch3:sec:projector}

The circuit of Sato \emph{et al.}~\cite{sato2024hamiltonian}
exponentiates ladder-string terms (alphabet $\{I,\sigma_{01},
\sigma_{10}\}$) via a Bell-basis change followed by one multi-controlled
$R_z$; Hu, Jin, Liu, and Zhang~\cite{hu2024quantum} extended the
framework to non-unitary PDE dynamics via Schr\"odingerisation, and Jin,
Liu, and Yu~\cite{jin2024heat} built gate-level circuits for the heat
equation with physical boundaries.  Relative to these we add the
\emph{projector} extension: strings containing
$z=\lvert0\rangle\langle0\rvert$ and $o=\lvert1\rangle\langle1\rvert$
(block projectors of separated multi-field systems) enter as plain
$0$-/$1$-controls; the Bell rotation outside the projected subspace is
undone exactly, so the circuit implements
$e^{-i\,\delta t\,(cP+\mathrm{h.c.})}$ for projector-valued $P$.
\ref{ch3:fig:circuits}(a) shows the compiled circuit for a
representative projector-containing term of the separated 2D wave
Hamiltonian, and \ref{ch3:fig:circuits}(b) the full Schr\"odingerised
pipeline it slots into.

\begin{figure}[tb]
  \centering
  \makebox[0pt][r]{(a)\hspace{0.96\linewidth}}\\[-0.5ex]
  \resizebox{\linewidth}{!}{%
  \begin{quantikz}[row sep={0.62cm,between origins}, column sep=0.22cm]
    \lstick{$q_0\;(\sigma_{10})$} & \qw & \targ{} & \qw & \qw
      & \ctrl{1} & \qw & \qw & \targ{} & \qw & \qw \\
    \lstick{$q_1\;(\sigma_{01})$} & \targ{} & \qw & \qw & \qw
      & \octrl{3} & \qw & \qw & \qw & \targ{} & \qw \\
    \lstick{$q_2\;(I)$} & \qw & \qw & \qw & \qw
      & \qw & \qw & \qw & \qw & \qw & \qw \\
    \lstick{$q_3\;(I)$} & \qw & \qw & \qw & \qw
      & \qw & \qw & \qw & \qw & \qw & \qw \\
    \lstick{$q_4\;(\sigma_{00})$} & \qw & \qw & \qw & \qw
      & \octrl{1} & \qw & \qw & \qw & \qw & \qw \\
    \lstick{$q_5\;(\sigma_{01},\,\mathrm{piv.})$}
      & \ctrl{-4} & \ctrl{-5} & \gate{P(\lambda)} & \gate{H}
      & \gate{R_z(2\gamma\,\delta t)} & \gate{H} & \gate{P(-\lambda)}
      & \ctrl{-5} & \ctrl{-4} & \qw
  \end{quantikz}}\\[1.2ex]
  \makebox[0pt][r]{(b)\hspace{0.96\linewidth}}\\[-0.5ex]
  \resizebox{\linewidth}{!}{%
  \begin{quantikz}[row sep={0.95cm,between origins}, column sep=0.18cm]
    \lstick{sys} & \qwbundle{n_s} & \gate{\lvert\psi_0\rangle} & \qw
      & \gate[2]{e^{+i\frac{\Delta t}{2}\eta\Sigma}}
      & \gate{e^{-i\Delta t H}}
      & \gate[2]{e^{+i\Delta t\,\eta\Sigma}}
      & \ \ldots\
      & \gate{e^{-i\Delta t H}}
      & \gate[2]{e^{+i\frac{\Delta t}{2}\eta\Sigma}}
      & \qw & \qw \\
    \lstick{$p$} & \qwbundle{n_p} & \gate{\lvert g\rangle}
      & \gate{\mathrm{QFT}^{\dagger}} & {} & \qw & {}
      & \ \ldots\
      & \qw & {} & \gate{\mathrm{QFT}} & \meter{p^{*}}
  \end{quantikz}}
  \caption{Explicit circuits.  (a)~Bell-basis evolution
  $e^{-i\,\delta t\,(cP+\mathrm{h.c.})}$ of the projector-containing
  term $P=\sigma_{01}\otimes\sigma_{00}\otimes I\otimes I\otimes
  \sigma_{01}\otimes\sigma_{10}$ of the separated 2D wave Hamiltonian
  (string \texttt{mzIImp}, qubit $q_k$ carrying string character $5-k$;
  $\lambda=\arg c$, $\gamma=|c|$; $\lambda=0$ for the real wave-equation
  coefficients): \textsc{cnot}s rotate the ladder qubits to the Bell
  basis, and the single multi-controlled $R_z$ carries the block
  projector $\sigma_{00}$ ($q_4$) as an \emph{open-circle}
  $0$-control---the entire compilation rule of
  Section~\protect\ref{ch3:sec:projector}.  The gate sequence is unitarily identical to
  the library-compiled circuit (verified to machine precision).
  (b)~The Schr\"odingerised pipeline (Section~\protect\ref{ch3:sec:circuit}) on bundled
  registers: product-state preparation
  $\lvert g\rangle\otimes\lvert\psi_0\rangle$, $\mathrm{QFT}^{\dagger}$
  on the $p$ register, Strang-split steps alternating the Bell wave
  circuit $e^{-i\Delta t H}$ with the diagonal damping phases
  $e^{+i\Delta t\,\eta\Sigma}$ (half steps at the ends), closing
  $\mathrm{QFT}$, and recovery by measuring the $p$ register,
  post-selecting the slice at $p^{*}$, and rescaling classically by
  $e^{p^{*}}$.}
  \label{ch3:fig:circuits}
\end{figure}

Due diligence (full-text check of both papers): neither
Ref.~\cite{hu2024quantum} nor Ref.~\cite{jin2024heat} circuitises
projector-containing strings---both use only the
$\{I,\sigma_{01},\sigma_{10}\}$ alphabet.  We confirmed this at code
level in the published companion notebook of
Ref.~\cite{hu2024quantum}: its term circuit ($W_j$) implements the
contiguous ladder pattern with rotations controlled uniformly on all
lower qubits (an all-or-nothing $X$-conjugation flag serving the
periodic corner terms), with no per-qubit $0$-/$1$-control selection
(the nilpotent shift needs no
projector guards in scalar single-field problems, and
Ref.~\cite{jin2024heat} routes boundary data into source terms precisely
to avoid non-Toeplitz structure).  Ref.~\cite{sato2025nonconservative}
comes closest: its symbolic operator algebra over
$\{I,\sigma_{00},\sigma_{11},\sigma_{01},\sigma_{10}\}$ produces
projector factors from boundary-truncated shifts and piecewise-constant
coefficients, and poly-depth Trotterisation is asserted by citation to
Ref.~\cite{sato2024hamiltonian}; but no gate-level compilation rule for
the projector factors and no projector-term circuit is given.  A caveat
we state explicitly: the
abstract class $\lvert a\rangle\langle b\rvert + \lvert b\rangle\langle
a\rvert$ over basis states in Lemma~1 of
Refs.~\cite{sato2024hamiltonian,hu2024quantum} implicitly \emph{contains}
Hermitian-paired projector strings as an existence statement; our
contribution is the explicit alphabet extension and compilation rule
(projector factor $\to$ $0$-/$1$-control) and its application to
multi-field block-coupled systems, not a new diagonalisation principle.
The extension is backward compatible.  Unitary-level validation: the
error is pure per-step Trotter error $0.50\,\delta t^{2}$ with no phase
offset.  This unlocks the separated 2D wave Hamiltonian
(\texttt{zm}/\texttt{mz} block couplings) needed for per-direction
absorbing layers.

\subsubsection{The Schr\"odingerised circuit}
\label{ch3:sec:circuit}

The full pipeline is
\begin{equation}
\begin{split}
  \mathrm{init}(g\otimes\psi_0)
  \;\to\; \mathrm{QFT}^{\dagger}_p
  \;\to\; [\text{Trotter steps}] \\
  \;\to\; \mathrm{QFT}_p
  \;\to\; \text{slice at } p^{*}.
\end{split}
  \label{ch3:eq:pipeline}
\end{equation}
Each step (diagonal $H_1$) factorises into the wave circuit (the Bell
construction of Section~\ref{ch3:sec:projector} on the system register) and a
diagonal damping phase with entries $e^{+i\,\delta t\,\eta_k\sigma_j}$.
Both Trotter levels are symmetrised (Strang outer; palindromic term order
inner), with measured clean second order: errors
$4\times10^{-2}\to10^{-2}\to2.7\times10^{-3}$ at $50/100/200$ steps on 17
qubits.

\emph{Fitted orders} [1D benchmark, error versus the exact
Schr\"odingerised evolution to isolate the splitting;
\ref{ch3:fig:convergence}(d)]: slope $-1.02$ (order~1), $-2.01$ (order~2
with both levels symmetrised); against the full matrix-exponential
reference the order-2 error saturates at the $p$-discretisation floor, as
expected.  \emph{Warping-profile order} [recovery error versus $n_p$,
fixed 1D PML case; \ref{ch3:fig:convergence}(c)]: $e^{-|p|}$ slope
${\sim}{-1}$ in $\Delta p$ ($4.1\times10^{-2}\to6.3\times10^{-4}$ over
$n_p=6$--$11$); $C^{1}$ cubic ${\sim}{-2}$
($1.3\times10^{-2}\to2.4\times10^{-6}$), confirming Remark~4.9 of
Ref.~\cite{jin2024recovery} in this setting.

\subsubsection{End-to-end demonstrations}
\label{ch3:sec:demos}

\begin{itemize}
\item 1D CPML, 14 qubits (6 system $+$ 8 $p$), 60 steps, $T=30$:
  recovered solution within $7\times10^{-4}$ of the exact non-unitary
  evolution.
\item 2D sponge, $8\times8$ grid, 17 qubits, $T=10$:
  $2.7\times10^{-3}$ at 200 steps (\ref{ch3:fig:gate17}).
\item 2D sponge, $32\times32$ grid, 21 qubits (12 system $+$ 9 $p$), up
  to 240 steps (983 native operations per step): errors
  $7.7\times10^{-3}$ / $1.7\times10^{-2}$ / $1.4\times10^{-2}$ at
  $T=8/16/24$; energy decays from $0.97$ to $1.25\times10^{-3}$
  (absorption complete); nine minutes of statevector simulation
  (\ref{ch3:fig:gate21}).
\end{itemize}
Qubit accounting: the absorbing overhead is $+1$ (sponge) or $+2$ (CPML)
block qubits plus $n_p\approx8$--$10$ for the $p$ register; $n_p$ depends
only logarithmically on $T$, the damping rate, and the target accuracy.

\begin{figure}
  \includegraphics[width=\linewidth]{chapter3/fig5_2d_gate_level_simulation.png}
  \caption{Gate-level absorbing 2D wave evolution (sponge layer,
  $8\times8$ grid, 17 qubits $=$ 6 grid $+$ 2 block $+$ 9 $p$-register).
  Top row: snapshots of the recovered field from the Trotterised quantum
  circuit at $t=2.5,5,7.5,10$.  Bottom row: the exact non-unitary
  reference $e^{At}\psi_0$.  The expanding pulse is absorbed in the
  boundary layer with no visible reflection; snapshot errors are at the
  $10^{-3}$ level (Section~\protect\ref{ch3:sec:demos}).}
  \label{ch3:fig:gate17}
\end{figure}

\begin{figure}
  \includegraphics[width=\linewidth]{chapter3/fig6_2d_gate_level_32x32.png}
  \caption{Gate-level absorbing 2D wave evolution at scale: $32\times32$
  grid, 21 qubits (12 system $+$ 9 $p$-register), up to 240 second-order
  Trotter steps (983 native operations per step).  Top row: recovered
  field from the quantum circuit at $t=0,8,16,24$ (dashed lines mark the
  absorbing-layer strips).  Bottom row: exact reference $e^{At}\psi_0$.
  The outgoing ring is absorbed across the layer; the total energy decays
  from $0.97$ to $1.25\times10^{-3}$, and the recovered field matches the
  reference to $7.7\times10^{-3}$--$1.7\times10^{-2}$.}
  \label{ch3:fig:gate21}
\end{figure}

% ======================================================================
\subsection{Complexity}
\label{ch3:sec:complexity}

Term counts inherit the QTT structure of the difference operators:
$2n_x+2$ terms in 1D and $4n_x+4$ in 2D ($n_x=n_y$), exactly linear in
the register size (fitted slopes $4.000/2.000$, intercepts $4/2$).
Transpiled to the $\{\mathrm{CX},u\}$ basis, one Strang wave step costs
$217/459/834/1343$ operations for $n_{\mathrm{sys}}=6/8/10/12$; a
power-law fit gives exponent $2.63$ with $r^{2}=0.99997$
(\ref{ch3:fig:gatecounts})---convincingly polynomial, and in fact
asymptotically quadratic: the term count is linear in $n$ and the
ancilla-free V-chain multi-controlled-$R_z$ is linear in its control
count, so ops/step $=O(n^{2})$.  The constants are fixed by the
measurements.  The per-step cost in the
$\{\textsc{cx}, U\}$ basis is reproduced to $0.03\%$
($r^{2}=0.9999999$) by the quadratic
\begin{equation}
  G_{\mathrm{step}}(n) \;=\; 16.7\,n^{2} - 112.7\,n + 293,
  \label{ch3:eq:stepcost}
\end{equation}
whose structure is transparent: a Strang step exponentiates $2n$
Hermitian-paired terms ($4n_x+4$ at $n=2n_x+2$), each one Bell basis
change plus one multi-controlled $R_z$ whose ancilla-free V-chain cost
is \emph{exactly} $32k-31$ elementary operations for $k\ge8$ controls
(measured), linear in $k$; with the QTT ladder structure averaging
$\bar k\approx n/4$ controls per term, the leading coefficient is
$2n\cdot32\,\bar k/n^{2}\approx16.7$.  The total gate count is steps
$\times$ ops/step with steps $= T/\delta t$ and $\delta t$ set by the
second-order bound $\delta t = \sqrt{\varepsilon_{\mathrm{tr}}/(C_2
T)}$, where $C_2$ is the Strang error constant of the slope-$(-2)$ fits
of \ref{ch3:fig:convergence}(d):
\begin{equation}
  G_{\mathrm{total}}
  \;=\; \sqrt{\frac{C_2\,T^{3}}{\varepsilon_{\mathrm{tr}}}}\;
        G_{\mathrm{step}}(n)
  \;=\; O\!\left(\frac{T^{3/2}}{\sqrt{\varepsilon_{\mathrm{tr}}}}\;
        n^{2}\right),
  \label{ch3:eq:totalcost}
\end{equation}
plus the damping diagonal per step and two $\mathrm{QFT}(n_p)$; the
power-law fit across the measured range reads $n^{2.63}$ because the
subleading terms of Eq.~\eqref{ch3:eq:stepcost} have not yet decayed at
$n\le12$.

\emph{Caveat (honest).}  The damping factor is a single native
\texttt{diagonal} instruction on the simulator, but a \emph{generic}
diagonal on $n$ qubits synthesises to $\Theta(2^{n})$ elementary gates:
we measure a CX count of exactly $2^{n}-2$ and totals within 15\% of the
$2^{n+1}-3$ synthesis bound.  For hardware the structure must be
exploited: the diagonal is $D_\eta\otimes\Sigma$ with $\eta$ linear in
the $p$ index and $\sigma$ a degree-$m$ polynomial profile on the layer
strips, so phase-polynomial constructions give $\mathrm{poly}(n)$
circuits, and the logic-minimisation synthesis of piecewise-constant
diagonal operators of Ref.~\cite{sato2025nonconservative}
(Espresso-compressed projector strings with matrix-product-state
coefficient oracles) provides an existing $\mathrm{poly}(n)$ toolbox
directly applicable to the strip-wise polynomial $\sigma$ profiles;
making this explicit is deferred to future work.

\begin{figure}
  \includegraphics[width=\linewidth]{chapter3/paper_P3_gatecounts.png}
  \caption{Gate cost per Trotter step after transpilation to the
  $\{\mathrm{CX},u\}$ basis, versus the number of system qubits $n$: 1D
  steps (squares) and 2D Strang steps (circles).  The dotted line is the
  power-law fit $\propto n^{2.63}$ ($r^{2}=0.99997$) to the 2D series;
  the term count itself is exactly linear in $n$ (Section~\protect\ref{ch3:sec:complexity}).}
  \label{ch3:fig:gatecounts}
\end{figure}

% ======================================================================
\subsection{Discussion and outlook}
\label{ch3:sec:discussion}

The $\lambda^{+}$ obstruction of Lemma~\ref{ch3:lem:threshold} is generic for
auxiliary-field absorbers---memory variables, auxiliary-differential-equation
formulations, split fields---because the injection/source asymmetry is
structural rather than an artefact of one discretisation.  The Lyapunov
symmetrizer is the corresponding general remedy, and its
quantum-implementable variants (block-encoding $S$, structured or banded
approximate symmetrizers) are open.  Natural next steps include a
hardware experiment in the style of Ref.~\cite{sato2024hamiltonian};
quantum-signal-processing or qubitisation in place of Trotter (and the
query-optimal Schr\"odingerisation formulation of
Ref.~\cite{jin2025queries}); CPML for Maxwell's equations
\cite{jin2023maxwell,ma2024maxwellcircuits}; and smoother warping
profiles (error-function or numerically optimal) beyond the $C^{1}$ cubic.

\medskip
\noindent\emph{Code and data availability.}  The implementation
(\texttt{lib/absorbing.py}, \texttt{lib/schrodingerisation.py}), all
experiment scripts and JSON result files
(\texttt{paper/experiments/}), and the executed demonstration notebook
are available in the accompanying repository at
\url{https://github.com/x-repos/qawave-cpml}.

% ======================================================================


\subsection{Proof of Lemma~\protect\ref{ch3:lem:threshold}}
\label{ch3:app:proof}

Fix a layer node $j$ with $\sigma_j>0$ and a stencil neighbour $i$ such
that the difference entry satisfies $|D_{ji}|=1/h$ (wave speed $c=1$;
the general case restores a factor $c$ in Eq.~\eqref{ch3:eq:lemmarate}).
Write $\beta=\sigma_j/\kappa_j+\alpha_j\ge0$ for the CFS decay rate,
$\gamma$ for the injection coefficient, and
$s=\sigma_j|D_{ji}|/(\kappa_j^{2}\gamma)$ for the memory-source
magnitude.  Let $G=\diag(g)>0$ be \emph{any} positive diagonal
similarity, and set
\begin{equation}
  u=\frac{g_{v_j}}{g_{\phi_j}}, \qquad
  t=\frac{g_{\phi_j}}{g_{w_i}}, \qquad
  \rho=u\,t=\frac{g_{v_j}}{g_{w_i}}.
  \label{ch3:eq:scalings}
\end{equation}
The transformed generator $\At = GAG^{-1}$ has unchanged diagonal, and
the three relevant coupling families scale as follows.
\begin{enumerate}
\item[(i)] \emph{Injection pair} $(v_j,\phi_j)$: entry $u\gamma$, reverse
  entry $0$, so $\Ht$ contains the principal $2\times2$ submatrix
  \begin{equation}
    \begin{bmatrix} 0 & u\gamma/2 \\ \bar u\bar\gamma/2 & -\beta
    \end{bmatrix}
    \;\Longrightarrow\;
    \lamax(\Ht) \ge f(|u|\gamma,\beta),
  \end{equation}
  with $f(x,\beta):=\bigl(\sqrt{\beta^{2}+x^{2}}-\beta\bigr)/2$.
\item[(ii)] \emph{Memory-source pair} $(w_i,\phi_j)$: entry $t\,s$
  analogously, giving $\lamax(\Ht)\ge f(|t|s,\beta)$.
\item[(iii)] \emph{Physical wave pair} $(v_j,w_i)$: elementwise,
  $\At[v,w]=-iD_{ji}\,\rho$ and $\At[w,v]=+iD_{ji}/\rho$ (using that $D$
  is real and $D^{\dagger}=-D_{\mathrm{b}}$), so
  \begin{equation}
  \begin{split}
    \Ht[v_j,w_i] &= \frac{iD_{ji}}{2}\bigl(\rho^{-1}-\rho\bigr) \\
    \Longrightarrow\quad
    \lamax(\Ht) &\ge \frac{|D_{ji}|}{2}\,\bigl|\rho^{-1}-\rho\bigr|.
  \end{split}
  \end{equation}
\end{enumerate}
All three are lower bounds on $\lamax(\Ht)$ by Cauchy interlacing (the
eigenvalues of any principal submatrix interlace those of the full
Hermitian matrix).

\emph{Case analysis.}  If $|\rho|\le1/2$ then
$|\rho^{-1}-\rho|\ge3/2$, so (iii) gives $\lamax(\Ht)\ge3/(4h)$.
Otherwise $|\rho|>1/2$ implies
\begin{equation}
  |u\gamma|\cdot|t\,s| \;=\; |\rho|\,\gamma s \;>\; \frac{\gamma s}{2}
  \;=\; \frac{\sigma_j|D_{ji}|}{2\kappa_j^{2}},
\end{equation}
hence
$\max(|u|\gamma,\,|t|s)\ge\sqrt{\sigma_j/(2h\kappa_j^{2})}$ and
(i)--(ii) give
$\lamax(\Ht)\ge f\bigl(\sqrt{\sigma_j/(2h\kappa_j^{2})},\,\beta\bigr)$.
Therefore, for every positive diagonal scaling,
\begin{equation}
  \lamax(\Ht) \;\ge\; \min\left\{\frac{3}{4h},\;
  \frac{\sqrt{\beta^{2}+\sigma_j/(2h\kappa_j^{2})}-\beta}{2}\right\},
  \label{ch3:eq:appbound}
\end{equation}
which is $\Theta\bigl(\sqrt{\sigma_j/h}\bigr)$ as $h\to0$ (the second
branch is the minimum once $\sigma_j/(2h)\gg\beta^{2}$, i.e.\ for any
fixed profile and small $h$).  The bound is independent of $\gamma$ and
of the scaling $G$.  The explicit balanced choice
$\gamma^{*}=\sqrt{2\sigma_{\max}/h}$ attains the same rate from above:
measured $\lambda^{+}=0.68$--$0.97$ across our configurations, versus the
lower bound $0.18$ at the $8\times8$, $\sigma_{\max}=0.5$ parameters---the
bound is a rate, not a sharp constant. \hfill$\blacksquare$

\medskip
\emph{Consistency check.}  The collapsed 1D form
(Section~\ref{ch3:sec:1d}) has no injection pair: the memory is eliminated and
the damping is diagonal, so the lemma does not apply---and indeed
$H_1=-\diag(\sigma)\preceq0$ there.
